\documentclass{readex-vivid}

\renewcommand{\readextitle}{讲义标题}
\renewcommand{\readexshorttitle}{讲义标题}
\renewcommand{\readexsubtitle}{结合多种资料重构的学习材料}
\renewcommand{\readexauthor}{Readex}
\renewcommand{\readexdate}{\today}
% \renewcommand{\readexcoverimage}{assets/cover.png}

% 默认启用 ReadOS 来源跳转。只有用户明确要求制作不含 ReadOS 超链接的
% 便携 PDF 时，才使用 \readexlinksfalse。
\readexlinkstrue

\begin{document}
\makereadexcover
\tableofcontents

\chapter{开始学习}

\section{为什么要学习这个主题}
先建立问题、动机和直觉，再进入形式化内容。

\begin{knowledgebox}
背景知识和帮助理解的类比放在这里。
\end{knowledgebox}

\section{核心思想与机制}
解释核心思想、工作机制、证据和例子。

\begin{importantbox}
用一句准确的话压缩本节最重要的结论。
\end{importantbox}

\chapter{总结与延伸}
总结核心关系、实际意义、开放问题和下一步。

\end{document}
